% generated by roam.py handbook export -- не редактировать, править заметки
% 20260923124012-спектральная_теорема_для_симметричных_матриц.org (body 3c7c31069135)

\begin{theorem}[{Спектральная теорема для симметричных матриц}]\label[theorem]{note:3BE0BBCE-AAB0-4CB4-8FA6-45277C0889CF}
  Пусть $A \in \R^{n \times n}$~--- симметричная матрица ($A = A^T$).
  Тогда в $\R^n$ существует ортонормированный базис из \hyperref[note:57D369D6-D748-41A2-9BDA-38D2FFC6601A]{собственных векторов} $A$:
  \[
  A q_i = \lambda_i q_i,\quad \left\langle q_i; q_j \right\rangle = \begin{cases} 1, & i = j \\ 0, & i \ne j \end{cases},\quad \lambda_i \in \R
  \]
\end{theorem}

\begin{remark}
  Матричная форма: $A = Q \Lambda Q^T$, где $Q = (q_1 \mid \dots \mid q_n)$~--- ортогональная матрица ($Q^T Q = I$), $\Lambda = \mathrm{diag}(\lambda_1, \dots, \lambda_n)$.
\end{remark}

\begin{remark}
  Главное следствие для оптимизации: в этом базисе квадратичная форма становится суммой квадратов. Если $h = \sum\limits_{i=1}^n c_i q_i$, то $\norm{h}^2 = \sum\limits_{i=1}^n c_i^2$ и
  \[
  \left\langle Ah; h \right\rangle = \sum\limits_{i=1}^n \lambda_i c_i^2, \qquad \lambda_{\min} \norm{h}^2 \le \left\langle Ah; h \right\rangle \le \lambda_{\max} \norm{h}^2
  \]
  Отсюда $A \ge \mu I \Leftrightarrow \lambda_{\min}(A) \ge \mu$ (см. \hyperref[note:94DADA28-87E4-4E36-B289-8415BCE750FC]{критерий сильной выпуклости в терминах вторых производных}).
\end{remark}

\begin{proof}
  Индукция по $n$. База $n = 1$: подходит $q_1 = 1$.

  \proofpart{Первый собственный вектор}
  По теореме о \hyperref[note:4F4E5575-ABE6-4D9B-9C61-B76A3AA9701C]{собственных значениях симметричных матриц} все \hyperref[note:9D2905FC-8E68-4E6F-9243-364269F795B0]{собственные значения} $A$ вещественны. Возьмем одно из них, $\lambda_1$.
  Матрица $A - \lambda_1 I$ вещественная и вырожденная, поэтому у нее есть ненулевое вещественное решение $(A - \lambda_1 I) q = 0$. Положим $q_1 = \frac{q}{\norm{q}}$.

  \proofpart{Ортогональное дополнение инвариантно}
  Пусть $W = \left\{ h \in \R^n : \left\langle h; q_1 \right\rangle = 0 \right\}$, $\dim W = n - 1$. Для $h \in W$ по симметричности $A$:
  \[
  \left\langle Ah; q_1 \right\rangle = h^T A^T q_1 = h^T A q_1 = \lambda_1 \left\langle h; q_1 \right\rangle = 0
  \]
  то есть $A W \subset W$.

  \proofpart{Переход}
  Выберем в $W$ ортонормированный базис $u_1, \dots, u_{n-1}$ и составим из него матрицу $U = (u_1 \mid \dots \mid u_{n-1})$ размера $n \times (n-1)$, $U^T U = I$.
  Для $w \in W$ верно $w = U U^T w$, а так как $A U$ состоит из векторов $W$, то $A U = U B$, где $B = U^T A U$.
  Матрица $B$ размера $(n-1) \times (n-1)$ симметрична: $B^T = U^T A^T U = B$. По предположению индукции у нее есть ортонормированный базис из собственных векторов $p_2, \dots, p_n$, $B p_i = \lambda_i p_i$.

  Положим $q_i = U p_i$ при $i = 2, \dots, n$. Тогда
  \[
  A q_i = A U p_i = U B p_i = \lambda_i U p_i = \lambda_i q_i, \qquad \left\langle q_i; q_j \right\rangle = p_i^T U^T U p_j = \left\langle p_i; p_j \right\rangle
  \]
  то есть $q_2, \dots, q_n$~--- ортонормированные собственные векторы $A$, лежащие в $W$, и потому ортогональные $q_1$.
  Вместе с $q_1$ получаем $n$ ортонормированных векторов~--- базис $\R^n$.
\end{proof}
